% chapters/05_performance.tex
\chapter{性能设计与评测}

\section{评测原则}

性能评测同时覆盖 Linux QEMU virtio 上的 virtio-blk 对照和 Asterinas KVM 中的保护 mapper。
两类环境分别报告，不使用跨平台结果计算倍率。每个有效工况固定运行五轮，保留原始
JSON、CSV、串口日志和环境元数据，并报告均值、中位数和标准差。

FIO 使用的诚实字段如下：

\begin{itemize}
    \item \texttt{iodepth\_set} 仅表示 fio 配置值；
    \item 无法观测真实并发时，\texttt{max\_\allowbreak in\_\allowbreak flight\_\allowbreak observed} 记为 \texttt{unavailable}；
    \item dm-verity 是只读目标，写入工况记为 N/A；
    \item Filebench 使用工具原始 MB/s，FIO 使用 MiB/s，不混用单位。
\end{itemize}

正式合同见仓库 \texttt{docs} 目录，最终结果见 \texttt{docs/\allowbreak evidence} 下的
\texttt{final-\allowbreak performance-\allowbreak 20260811} 目录。

\section{真实性审计}

早期 Asterinas dm-crypt FIO 在 mapper 初始化失败后继续执行，fio 将不存在的设备
路径创建为普通文件，产生了数量级异常。该批结果被分类为 \texttt{HARNESS\_\allowbreak BUG}，
未进入正式矩阵，旧数字和失败重试也不作为生成器输入。重新执行的 dm-crypt 使用
guest 内真实 \texttt{/dev/\allowbreak mapper/\allowbreak dm-\allowbreak crypt-\allowbreak bench}，dm-verity 则使用修正后的真实
mapper runner。

\section{FIO 结果}

Asterinas 在 256 MiB 正式合同下给出 raw / dm-crypt / dm-verity 配对保留率（4 KiB 同步读、QD1、五轮中位）。raw 行合并两组共十次测量；相对值使用同一镜像的配对 raw，不与后文 13 格 mapper 矩阵混算。

\begin{table}[H]
    \centering
    \caption{Asterinas 256 MiB 正式合同（MiB/s）}
    \begin{tabular}{p{0.28\textwidth}p{0.22\textwidth}p{0.22\textwidth}}
        \toprule
        \textbf{设备} & \textbf{中位吞吐} & \textbf{相对配对 raw} \\
        \midrule
        raw（两组共十轮） & 4.249 & 100.0\% \\
        dm-crypt & 4.059 & 95.6\% \\
        dm-verity & 3.815 & 87.3\% \\
        \bottomrule
    \end{tabular}
\end{table}

同一宿主机、KVM、CPU=\texttt{host}、1 vCPU、8\,GiB、同一 virtio-blk 数据盘前 256\,MiB、4\,KiB、\texttt{sync}、\texttt{direct=1}、QD1 的 Linux guest 对照如下。dm-crypt 保留率相对表中 raw 列；dm-verity 保留率相对配对 raw（Linux 4.641、Asterinas 4.368）。Asterinas 相对 Linux guest 的绝对吞吐为 raw 82.4\%、dm-crypt 85.5\%、dm-verity 87.3\%（$3.815/4.369$）。自身保留率不代表绝对吞吐超过 Linux。

\begin{table}[H]
    \centering
    \small
    \caption{Linux guest 与 Asterinas 同条件对照（MiB/s）}
    \begin{tabular}{lrrrrr}
        \toprule
        \textbf{系统} & \textbf{raw} & \textbf{dm-crypt} & \textbf{crypt 保留率} & \textbf{dm-verity} & \textbf{verity 保留率} \\
        \midrule
        Linux guest & 5.150 & 4.746 & 92.158\% & 4.369 & 94.1\% \\
        Asterinas & 4.245 & 4.059 & 95.621\% & 3.815 & 87.3\% \\
        \bottomrule
    \end{tabular}
\end{table}

Linux QEMU virtio 对照在 virtio-blk 上测试 raw、dm-crypt 与 dm-verity 读，
统一使用 \texttt{libaio}、\texttt{direct=1}、\texttt{numjobs=1}。代表性五轮中位数如下：

\begin{table}[H]
    \centering
    \caption{Linux QEMU virtio FIO 五轮中位带宽（MiB/s）}
    \begin{tabular}{p{0.25\textwidth}p{0.15\textwidth}p{0.15\textwidth}p{0.18\textwidth}p{0.18\textwidth}}
        \toprule
        \textbf{负载} & \textbf{块大小} & \textbf{QD} & \textbf{raw} & \textbf{dm-crypt} \\
        \midrule
        read & 4 KiB & 1 & 11.391 & 10.942 \\
        read & 256 KiB & 1 & 134.623 & 129.474 \\
        randread & 4 KiB & 16 & 58.931 & 58.947 \\
        randwrite & 256 KiB & 1 & 130.124 & 109.453 \\
        \bottomrule
    \end{tabular}
\end{table}

Asterinas KVM 真实保护 mapper 五轮中位：

\begin{table}[H]
    \centering
    \caption{Asterinas KVM FIO 五轮中位带宽（MiB/s）}
    \begin{tabular}{p{0.18\textwidth}p{0.12\textwidth}p{0.10\textwidth}p{0.16\textwidth}p{0.16\textwidth}p{0.16\textwidth}}
        \toprule
        \textbf{负载} & \textbf{块大小} & \textbf{QD} & \textbf{raw} & \textbf{dm-crypt} & \textbf{dm-verity} \\
        \midrule
        read & 4 KiB & 1 & 25.080 & 8.908 & 7.505 \\
        read & 256 KiB & 1 & 1230.577 & 97.456 & 53.363 \\
        randread & 4 KiB & 16 & 25.315 & 8.614 & 4.618 \\
        randwrite & 256 KiB & 1 & 1091.016 & 93.541 & N/A（只读） \\
        \bottomrule
    \end{tabular}
\end{table}

正式矩阵共有 Asterinas dm-crypt 13 项、dm-verity 7 项、Linux raw 13 项和 Linux
dm-crypt 13 项，每项五轮；另有 6 个 dm-verity 写入 N/A。

\section{Filebench、SQLite 与启动耗时}

Asterinas 与 Linux initramfs 的 Filebench / SQLite 对照在同一宿主机、QEMU/KVM、
\textbf{1 vCPU}、guest 内存 \textbf{4 GiB}、CPU 模型 \texttt{Icelake-Server,+x2apic}
（嵌套 VMware，不用 \texttt{-cpu host}）、\texttt{RELEASE\_LTO=1}、不开 mlock。
两边使用同一份 initramfs 作业和 ext2 virtio-blk。下表为单次实测，不是五轮中位。
Filebench 为 fileserver，\texttt{nfiles=10000}，\texttt{nthreads=10}，
\texttt{filesize=131072}，runtime 60\,s。SQLite 为官方作业
\verb|sqlite-speedtest1 --size 1000 /ext2|（TOTAL 含 VACUUM），Asterinas
使用上游默认 \texttt{SCHEME=iommu}。证据目录
\texttt{docs/evidence/asterinas-linux-initramfs-4g-20260818/}。

\begin{table}[H]
    \centering
    \caption{4 GiB Asterinas vs Linux initramfs（单次）}
    \begin{tabular}{p{0.42\textwidth}p{0.24\textwidth}p{0.24\textwidth}}
        \toprule
        \textbf{测试} & \textbf{Asterinas} & \textbf{Linux initramfs} \\
        \midrule
        Filebench fileserver（ops/s） & 8398.724 & 11545.120 \\
        Filebench fileserver（MB/s） & 201.9 & 278.5 \\
        SQLite \texttt{speedtest1} TOTAL（s） & 111.956 & 55.302 \\
        \bottomrule
    \end{tabular}
\end{table}

Linux QEMU guest 上 raw / dm-crypt 的 Filebench（8 GiB、\texttt{-cpu host}、
2000 文件、50 线程、五轮中位）与 10000 行 SQLite 是另一套口径，只用于看
dm-crypt 相对 raw 的应用层开销，不能与上表横比。dm-verity 与写密集
fileserver 不兼容，记为 N/A；SQLite dm-verity 只执行查询。启动时间从 QEMU
进程启动计至串口出现受保护根文件系统 PASS 标记；相对 plain 的中位开销为
dm-crypt 0.582\%，dm-verity 2.751\%。

\begin{table}[H]
    \centering
    \caption{Linux guest 应用层与 Asterinas 启动（五轮中位）}
    \begin{tabular}{p{0.38\textwidth}p{0.18\textwidth}p{0.18\textwidth}p{0.16\textwidth}}
        \toprule
        \textbf{测试} & \textbf{raw/plain} & \textbf{dm-crypt} & \textbf{dm-verity} \\
        \midrule
        Filebench fileserver（ops/s） & 27855.155 & 25666.016 & N/A \\
        Filebench fileserver（MB/s） & 656.9 & 605.0 & N/A \\
        SQLite 10000 行总耗时（s） & 0.138095 & 0.143336 & 0.001151（只读） \\
        Asterinas 启动耗时（s） & 51.331303 & 51.629820 & 52.743415 \\
        \bottomrule
    \end{tabular}
\end{table}

\section{实现侧性能设计}

dm-crypt 使用 AES-NI 后端和有界临时缓冲池，降低 AES-XTS 与频繁分配的开销；
dm-verity 复用已完成 BIO 数据，并对经过完整认证的 hash block 使用有界缓存。
缓存遵循“先认证、后插入”，命中不会降低 root hash 信任边界。所有优化仍以
R01--R03、C01--C05 和 V01--V07 正确性测试为前置门禁。

\section{小结}

最终报告只接受能够定位到真实 block device 或 mapper、参数完整且轮数满足合同的
数据。完整 46 个 PASS 单元格、6 个 N/A、逐轮统计和证据路径见
\texttt{PERFORMANCE\_\allowbreak MATRIX.csv} 与中文性能验证报告。